\subsection{QUASI-COMPACT SPACES AND COMPACT SPACES}

DEFINITION 1. A topological space X is said to be quasi-compact if it satisfies the following axiom:

\begin{enumerate}
\def\labelenumi{(\Alph{enumi})}
\setcounter{enumi}{2}
\tightlist
\item
  Every filter on X has at least one cluster point.
\end{enumerate}

A topological space is said to be compact if it is quasi-compact and Hausdorff.

It follows immediately from this axiom that if \(f\) is a mapping of a set \(Z\) into a quasi-compact space \(X\) , and \(\mathfrak{F}\) is any filter on \(Z\) , then \(f\) has at least one cluster point with respect to \(\mathfrak{F}\) . In particular, every sequence of points of a quasi-compact space has at least one cluster point; but this condition is not equivalent to (C) (Exercise 11).

We give three axioms each of which is equivalent to axiom (C):

(C') Every ultrafilter on \(X\) is convergent.

\((\mathrm{C}^{\prime}) \Longrightarrow (\mathrm{C})\) : If \(\mathfrak{F}\) is a filter on \(\mathbf{X}\) then there is an ultrafilter finer than \(\mathfrak{F}\) (§ 6, no. 4, Theorem 1). Since this ultrafilter converges to a point \(x\) , \(x\) is a cluster point of \(\mathfrak{F}\) .

\begin{enumerate}
\def\labelenumi{(\Alph{enumi})}
\setcounter{enumi}{2}
\tightlist
\item
  \(\Longrightarrow\) (C'): For if an ultrafilter has a cluster point then it converges to this point (§ 7, no. 2, Corollary to Proposition 4).
\end{enumerate}

If \(f\) is a mapping of a set \(Z\) into a quasi-compact space \(X\) , and \(\mathfrak{U}\) is an ultrafilter on \(Z\) , \(f\) has at least one limit point with respect to \(\mathfrak{U}\) (§ 6, no. 6, Proposition 10).

(C'') Every family of closed subsets of X whose intersection is empty contains a finite subfamily whose intersection is empty.

\begin{enumerate}
\def\labelenumi{(\Alph{enumi})}
\setcounter{enumi}{2}
\tightlist
\item
  \(\Longrightarrow\) (C'\,'): Let \(\mathfrak{G}\) be a family of closed subsets of X with empty intersection. If every finite subfamily of \(\mathfrak{G}\) has a non-empty intersection, then \(\mathfrak{G}\) generates a filter (§ 6, no. 2, Proposition 1) which has a cluster point by hypothesis. This point belongs to all the sets of \(\mathfrak{G}\) (since they are closed); so we have a contradiction.
\end{enumerate}

\((\mathbf{C}^{\prime\prime}) \Longrightarrow (\mathbf{C}) : \text{For if } (\mathbf{C}) \text{ is false then there is a filter } \mathfrak{F} \text{ on } \mathbf{X} \text{ which has no cluster point; hence the closures of the sets of } \mathfrak{F} \text{ form a family of closed subsets of } \mathbf{X} \text{ contradicting axiom } (\mathbf{C}^{\prime\prime}).\)

(C \(^{iii}\) ) (Axiom of Borel-Lebesgue) Every open covering of X contains a finite open covering of X.

\((\mathbf{C}^{\prime\prime\prime}) \iff (\mathbf{C}^{\prime\prime})\) by taking complements.

If X is quasi-compact, then every locally finite covering R of X is finite. For each point of X has an open neighbourhood which meets only a finite number of sets of R, and by \((\mathbf{C}^{\prime\prime})\) a finite number of these neighbourhoods covers X.

Examples. 1) Every finite space is quasi-compact, and more generally every space in which there is only a finite number of open sets is quasi-compact. A finite space is compact if and only if it is discrete, for a finite Hausdorff space is discrete (§ 8, no.1, Corollary to Proposition 3). Conversely, every compact discrete space is finite, for in such a space the sets consisting of a single point are open; hence the space is finite by (C'\,').

\begin{enumerate}
\def\labelenumi{\arabic{enumi})}
\setcounter{enumi}{1}
\tightlist
\item
  Let X be a set, and give X the topology in which the closed sets are X and all finite subsets of X {[}this set of subsets clearly satisfies axioms \((\mathrm{O}_{\mathrm{I}}^{\prime})\) and \((\mathrm{O}_{\mathrm{II}}^{\prime})\) of §1, no. 4{]}. The topological space so defined is quasi-compact. For if \((\mathbf{F}_{i})_{i\in I}\) is a family of closed subsets of X with empty intersection, then \(F_{\alpha}\) is finite for some \(\alpha\in I\) . Let \(a_{k}\) ( \(1\leqslant k\leqslant n\) ) be the elements of \(F_{\alpha}\) ; then by hypothesis for each index k there is an index \(i_{k}\in I\) such that \(a_{k}\notin F_{i_{k}}\) ; the intersection of the \(F_{i_{k}}\) ( \(1\leqslant k\leqslant n\) ) and \(F_{\alpha}\) is therefore empty, whence axiom \((\mathrm{C}^{\prime\prime})\) is satisfied. If X is infinite it is not Hausdorff.
\end{enumerate}

Remark. Quasi-compact (non-Hausdorff) spaces are of use mainly in applications of topology to algebraic geometry and are seldom featured in other mathematical theories, where on the contrary compact spaces play an important role.

THEOREM 1. Let \(\mathfrak{F}\) be a filter on a quasi-compact space \(\mathbf{X}\) and let \(\mathbf{A}\) be the set of cluster points of \(\mathfrak{F}\) . Then every neighbourhood of \(\mathbf{A}\) belongs to \(\mathfrak{F}\) .

Let V be a neighbourhood of A and suppose it possible that every set of F meets CV. The intersections of the sets of F with CV then form a base of a filter G on X; X is quasi-compact, so that G has at least one cluster point y, which does not belong to A, because the neighbourhood V of A does not meet some of the sets of G. But since G is finer than F, y is also a cluster point of F, which is contrary to hypothesis.

COROLLARY. For a filter on a compact space to converge it is necessary and sufficient that it has a single cluster point.

Necessity by § 8, no. 1, Proposition 1; sufficiency by Theorem 1 above.

